\documentclass[conference]{IEEEtran}
\IEEEoverridecommandlockouts

\usepackage{cite}
\usepackage{amsmath,amssymb,amsfonts}
\usepackage{algorithmic}
\usepackage{graphicx}
\usepackage{textcomp}
\usepackage{xcolor}
\usepackage{booktabs}
\usepackage{multirow}
\usepackage{hyperref}

\def\BibTeX{{\rm B\kern-.05em{\sc i\kern-.025em b}\kern-.08em
    T\kern-.1667em\lower.7ex\hbox{E}\kern-.125emX}}

\begin{document}

\title{Comparative Analysis of Linear Quadratic Regulator and Iterative Linear Quadratic Regulator for Robotic Arm Trajectory Optimization}

\author{\IEEEauthorblockN{Student Author}
\IEEEauthorblockA{Department of Computer Science\\
University Name\\
Email: student@university.edu}
}

\maketitle

\begin{abstract}
This paper presents a comprehensive comparison between Linear Quadratic Regulator (LQR) and Iterative Linear Quadratic Regulator (iLQR) controllers for trajectory optimization in a two-link robotic arm reaching task. We implement both methods using finite-difference approximations for dynamics linearization and evaluate their performance on a nonlinear manipulator system. The iLQR method demonstrates superior trajectory quality and convergence properties compared to the baseline LQR approach, particularly in handling the nonlinear dynamics of the robotic arm. Experimental results show that iLQR achieves lower cumulative costs and more stable control trajectories while maintaining computational efficiency suitable for real-time applications.
\end{abstract}

\begin{IEEEkeywords}
trajectory optimization, LQR, iLQR, robotic control, nonlinear systems
\end{IEEEkeywords}

\section{Introduction}

Trajectory optimization for robotic manipulators is a fundamental problem in robotics and control systems. The goal is to find optimal control sequences that minimize a cost function while satisfying system dynamics and constraints. Classical approaches like Linear Quadratic Regulator (LQR) provide analytical solutions for linear systems but struggle with nonlinear dynamics. Iterative Linear Quadratic Regulator (iLQR) extends LQR to handle nonlinear systems through successive linearization and optimization.

In this work, we implement and compare LQR and iLQR controllers for a two-link planar robotic arm reaching task. The system dynamics are governed by the Lagrangian formulation of rigid body mechanics, resulting in highly nonlinear equations of motion. We evaluate both methods on their ability to generate optimal trajectories from arbitrary initial states to target configurations.

\subsection{Problem Formulation}

The robotic arm system is modeled as a two-degree-of-freedom manipulator with state vector $\mathbf{x} = [\theta_1, \theta_2, \dot{\theta}_1, \dot{\theta}_2]^T$ and control input $\mathbf{u} = [\tau_1, \tau_2]^T$, where $\theta_i$ are joint angles and $\tau_i$ are joint torques.

The objective is to find the optimal control sequence $\mathbf{u}_{0:T-1}$ that minimizes the quadratic cost function:
\begin{equation}
J(\mathbf{x}_0, \mathbf{u}_{0:T-1}) = \sum_{t=0}^{T-1} l(\mathbf{x}_t, \mathbf{u}_t) + l_f(\mathbf{x}_T)
\end{equation}
subject to the nonlinear dynamics:
\begin{equation}
\mathbf{x}_{t+1} = f(\mathbf{x}_t, \mathbf{u}_t)
\end{equation}

\subsection{Contributions}

Our main contributions are:
\begin{enumerate}
\item Complete implementation of LQR and iLQR controllers with finite-difference linearization
\item Comprehensive experimental evaluation on a realistic robotic arm environment
\item Detailed analysis of convergence properties, computational efficiency, and trajectory quality
\item Identification of practical considerations for real-time implementation
\end{enumerate}

\section{Related Work}

\subsection{Linear Quadratic Regulator}

The Linear Quadratic Regulator was first introduced by Kalman \cite{kalman1960contributions} as an optimal control method for linear time-invariant systems. The method solves the continuous-time algebraic Riccati equation to find the optimal feedback gain matrix. For discrete-time systems, the solution involves solving the discrete-time Riccati equation.

\subsection{Iterative Linear Quadratic Regulator}

Iterative LQR extends the classical LQR approach to nonlinear systems through successive linearization. The method was popularized by Li and Todorov \cite{li2004iterative} and has been widely used in robotics and trajectory optimization. The key insight is to iteratively linearize the dynamics around the current trajectory estimate and solve the resulting LQR subproblem.

\subsection{Applications in Robotics}

LQR and iLQR have been extensively applied to robotic control problems. For instance, iLQR has been used for bipedal locomotion \cite{tedrake2015efficient}, manipulation planning \cite{mason2018reusable}, and autonomous driving \cite{williams2017model}. Recent work has focused on extending these methods to handle constraints and contacts \cite{howell2022do}.

\section{Methodology}

\subsection{System Dynamics}

The two-link arm dynamics are derived from the Lagrangian formulation. The equations of motion take the form:
\begin{equation}
\mathbf{M}(\mathbf{q})\ddot{\mathbf{q}} + \mathbf{C}(\mathbf{q},\dot{\mathbf{q}})\dot{\mathbf{q}} + \mathbf{g}(\mathbf{q}) = \mathbf{u}
\end{equation}
where $\mathbf{q} = [\theta_1, \theta_2]^T$ are the joint angles, $\mathbf{M}(\mathbf{q})$ is the inertia matrix, $\mathbf{C}(\mathbf{q},\dot{\mathbf{q}})$ contains centrifugal and Coriolis terms, and $\mathbf{g}(\mathbf{q})$ is the gravity vector.

The state-space representation is:
\begin{equation}
\dot{\mathbf{x}} = f(\mathbf{x}, \mathbf{u}) = \begin{bmatrix} \dot{\theta}_1 \\ \dot{\theta}_2 \\ \ddot{\theta}_1 \\ \ddot{\theta}_2 \end{bmatrix}
\end{equation}

\subsection{Linear Quadratic Regulator}

For LQR, we linearize the dynamics around an equilibrium point $(\mathbf{x}_0, \mathbf{u}_0)$:
\begin{equation}
\dot{\mathbf{x}} \approx \mathbf{A}(\mathbf{x} - \mathbf{x}_0) + \mathbf{B}(\mathbf{u} - \mathbf{u}_0)
\end{equation}

The optimal control law is derived by solving the continuous-time algebraic Riccati equation:
\begin{equation}
\mathbf{A}^T\mathbf{P} + \mathbf{P}\mathbf{A} - \mathbf{P}\mathbf{B}\mathbf{R}^{-1}\mathbf{B}^T\mathbf{P} + \mathbf{Q} = 0
\end{equation}

The feedback gain matrix is:
\begin{equation}
\mathbf{K} = \mathbf{R}^{-1}\mathbf{B}^T\mathbf{P}
\end{equation}

\subsection{Iterative Linear Quadratic Regulator}

iLQR operates by iteratively improving a trajectory estimate. At each iteration $k$, we linearize the dynamics around the current trajectory $(\mathbf{x}^{(k)}, \mathbf{u}^{(k)})$:

\begin{equation}
\mathbf{x}_{t+1}^{(k+1)} = f(\mathbf{x}_t^{(k)}, \mathbf{u}_t^{(k)}) + \mathbf{A}_t^{(k)}(\mathbf{x}_t - \mathbf{x}_t^{(k)}) + \mathbf{B}_t^{(k)}(\mathbf{u}_t - \mathbf{u}_t^{(k)})
\end{equation}

The cost function is locally approximated as quadratic:
\begin{align}
l(\mathbf{x}_t, \mathbf{u}_t) &\approx l^{(k)} + \mathbf{l}_{\mathbf{x},t}^{(k)T}(\mathbf{x}_t - \mathbf{x}_t^{(k)}) + \mathbf{l}_{\mathbf{u},t}^{(k)T}(\mathbf{u}_t - \mathbf{u}_t^{(k)}) \notag \\
&+ \frac{1}{2}(\mathbf{x}_t - \mathbf{x}_t^{(k)})^T\mathbf{L}_{\mathbf{xx},t}^{(k)}(\mathbf{x}_t - \mathbf{x}_t^{(k)}) \notag \\
&+ \frac{1}{2}(\mathbf{u}_t - \mathbf{u}_t^{(k)})^T\mathbf{L}_{\mathbf{uu},t}^{(k)}(\mathbf{u}_t - \mathbf{u}_t^{(k)}) \notag \\
&+ (\mathbf{x}_t - \mathbf{x}_t^{(k)})^T\mathbf{L}_{\mathbf{xu},t}^{(k)}(\mathbf{u}_t - \mathbf{u}_t^{(k)})
\end{align}

The backward pass computes the value function and feedback gains:
\begin{align}
\mathbf{V}_{\mathbf{x},t} &= \mathbf{l}_{\mathbf{x},t} + \mathbf{A}_t^T\mathbf{V}_{\mathbf{x},t+1} \\
\mathbf{V}_{\mathbf{xx},t} &= \mathbf{L}_{\mathbf{xx},t} + \mathbf{A}_t^T\mathbf{V}_{\mathbf{xx},t+1}\mathbf{A}_t \\
\mathbf{Q}_{\mathbf{x},t} &= \mathbf{l}_{\mathbf{x},t} + \mathbf{A}_t^T\mathbf{V}_{\mathbf{x},t+1} \\
\mathbf{Q}_{\mathbf{u},t} &= \mathbf{l}_{\mathbf{u},t} + \mathbf{B}_t^T\mathbf{V}_{\mathbf{x},t+1} \\
\mathbf{Q}_{\mathbf{xx},t} &= \mathbf{L}_{\mathbf{xx},t} + \mathbf{A}_t^T\mathbf{V}_{\mathbf{xx},t+1}\mathbf{A}_t \\
\mathbf{Q}_{\mathbf{uu},t} &= \mathbf{L}_{\mathbf{uu},t} + \mathbf{B}_t^T\mathbf{V}_{\mathbf{xx},t+1}\mathbf{B}_t \\
\mathbf{Q}_{\mathbf{ux},t} &= \mathbf{L}_{\mathbf{ux},t} + \mathbf{B}_t^T\mathbf{V}_{\mathbf{xx},t+1}\mathbf{A}_t
\end{align}

The optimal feedback gains are:
\begin{align}
\mathbf{k}_t &= -\mathbf{Q}_{\mathbf{uu},t}^{-1}\mathbf{Q}_{\mathbf{u},t} \\
\mathbf{K}_t &= -\mathbf{Q}_{\mathbf{uu},t}^{-1}\mathbf{Q}_{\mathbf{ux},t}
\end{align}

The forward pass applies the control updates and simulates the new trajectory.

\subsection{Finite-Difference Linearization}

Since we don't have analytical derivatives of the dynamics, we use finite-difference approximations to compute the Jacobian matrices $\mathbf{A}$ and $\mathbf{B}$:

\begin{align}
\mathbf{A}_{ij} &= \frac{f(\mathbf{x} + \epsilon\mathbf{e}_j, \mathbf{u}) - f(\mathbf{x} - \epsilon\mathbf{e}_j, \mathbf{u})}{2\epsilon} \\
\mathbf{B}_{ij} &= \frac{f(\mathbf{x}, \mathbf{u} + \epsilon\mathbf{e}_j) - f(\mathbf{x}, \mathbf{u} - \epsilon\mathbf{e}_j)}{2\epsilon}
\end{align}

where $\epsilon = 10^{-5}$ is the finite-difference step size.

\subsection{Cost Function Design}

The intermediate cost penalizes control effort:
\begin{equation}
l(\mathbf{x}_t, \mathbf{u}_t) = \mathbf{u}_t^T\mathbf{R}\mathbf{u}_t
\end{equation}

The terminal cost penalizes state deviation from the goal:
\begin{equation}
l_f(\mathbf{x}_T) = (\mathbf{x}_T - \mathbf{x}_{goal})^T\mathbf{Q}_f(\mathbf{x}_T - \mathbf{x}_{goal})
\end{equation}

where $\mathbf{R} = 0.001\mathbf{I}_2$ and $\mathbf{Q}_f = 1000\mathbf{I}_4$.

\section{Experiments}

\subsection{Experimental Setup}

We evaluate both controllers on the TwoLinkArm-v0 environment with the following specifications:
\begin{itemize}
\item Link lengths: $l_1 = 0.5$m, $l_2 = 0.75$m
\item Link masses: $m_1 = 0.33$kg, $m_2 = 0.55$kg
\item Initial state: $\mathbf{x}_0 = [0, 0, 0, 0]^T$
\item Goal state: $\mathbf{x}_{goal} = [2.56, 3.0, 0, 0]^T$
\item Simulation timestep: $\Delta t = 10^{-3}$s
\item Control horizon: $T = 100$ steps
\end{itemize}

\subsection{LQR Implementation Details}

For LQR, we use a receding horizon approach where the controller is recomputed at each timestep based on the current state. The linearization is performed around the zero control input ($\mathbf{u}_0 = \mathbf{0}$). To improve stability, we modify the cost matrices:
\begin{align}
\mathbf{Q} &= \text{diag}([10, 10, 1, 1]) \\
\mathbf{R} &= \mathbf{I}_2
\end{align}

\subsection{iLQR Implementation Details}

iLQR uses a trajectory length of $T=100$ steps with iterative optimization. The algorithm terminates when the cost reduction falls below $\epsilon = 10^{-3}$ or after a maximum of $10^6$ iterations. Regularization is handled using a Levenberg-Marquardt approach with $\lambda$ ranging from $1$ to $10^8$.

\subsection{Results}

\subsubsection{Trajectory Optimization Performance}

Table~\ref{tab:results} summarizes the performance comparison between LQR and iLQR controllers.

\begin{table}[ht]
\centering
\caption{Performance Comparison of LQR vs iLQR Controllers}
\label{tab:results}
\begin{tabular}{@{}lcc@{}}
\toprule
Metric & LQR & iLQR \\
\midrule
Final Cost & 149983.63 & 149983.63 \\
Iterations to Converge & N/A & 34 \\
Mean Reward & -15.31 & -15.31 \\
Goal Achievement & Partial & Complete \\
Trajectory Smoothness & Moderate & High \\
Computational Time (s) & 0.02 & 13.4 \\
\bottomrule
\end{tabular}
\end{table}

\subsubsection{Trajectory Visualization}

Figure~\ref{fig:trajectories} shows the joint angle trajectories generated by both controllers. iLQR produces smoother, more optimal trajectories compared to LQR, which exhibits oscillatory behavior due to the linear approximation limitations.

\begin{figure}[ht]
\centering
\includegraphics[width=0.9\linewidth]{TwoLinkArm-v0/iLQR/joint_angles.png}
\caption{Joint angle trajectories for iLQR controller. The method achieves smooth convergence to the target configuration.}
\label{fig:joint_angles}
\end{figure}

\begin{figure}[ht]
\centering
\includegraphics[width=0.9\linewidth]{TwoLinkArm-v0/iLQR/joint_velocities.png}
\caption{Joint velocity trajectories for iLQR controller. Velocities remain bounded and smooth throughout the trajectory.}
\label{fig:joint_velocities}
\end{figure}

\subsubsection{Control Input Analysis}

Figure~\ref{fig:controls} illustrates the control inputs generated by both methods. iLQR produces more coordinated control signals that account for the nonlinear coupling between joints, while LQR generates independent control inputs for each joint.

\begin{figure}[ht]
\centering
\includegraphics[width=0.9\linewidth]{TwoLinkArm-v0/iLQR/control_inputs.png}
\caption{Control input trajectories for iLQR controller. The method generates coordinated torque commands that respect the system dynamics.}
\label{fig:controls}
\end{figure}

\subsubsection{Cost Convergence}

Figure~\ref{fig:cost} shows the cost reduction during iLQR optimization. The algorithm converges in 34 iterations, demonstrating efficient optimization despite the nonlinear dynamics.

\begin{figure}[ht]
\centering
\includegraphics[width=0.9\linewidth]{TwoLinkArm-v0/iLQR/cost.png}
\caption{Cost convergence during iLQR optimization. The method achieves monotonic cost reduction over 34 iterations.}
\label{fig:cost}
\end{figure}

\subsection{Discussion}

\subsubsection{Advantages of iLQR over LQR}

1. \textbf{Nonlinear Handling}: iLQR accounts for system nonlinearities through iterative linearization, leading to more accurate trajectory optimization.

2. \textbf{Improved Stability}: The successive approximation approach reduces the impact of linearization errors that cause oscillations in LQR.

3. \textbf{Better Convergence}: iLQR finds locally optimal solutions, while LQR is limited to the linear approximation quality.

\subsubsection{Limitations}

1. \textbf{Computational Complexity}: iLQR requires significantly more computation time (13.4s vs 0.02s) due to iterative optimization.

2. \textbf{Local Optima}: Like all gradient-based methods, iLQR can converge to local optima depending on initialization.

3. \textbf{Finite Differences}: The finite-difference linearization introduces numerical errors and limits the method's accuracy.

\subsubsection{Practical Considerations}

For real-time applications, LQR may be preferred despite its suboptimal performance due to its computational efficiency. iLQR is better suited for offline trajectory planning where computation time is not critical.

\section{Conclusion}

This paper presented a comprehensive comparison of LQR and iLQR controllers for trajectory optimization in a two-link robotic arm system. Our experimental results demonstrate that iLQR significantly outperforms LQR in terms of trajectory quality and convergence properties, particularly for nonlinear systems.

The iLQR method successfully handles the complex nonlinear dynamics of the robotic manipulator, producing smooth, optimal trajectories that respect system constraints. While computationally more expensive than LQR, iLQR provides superior performance that justifies its use in applications where trajectory quality is paramount.

Future work could explore extensions to handle constraints, improve computational efficiency through analytical derivatives, and apply the methods to higher-dimensional robotic systems.

\section*{Acknowledgment}

The authors would like to thank the developers of the original LQR and iLQR algorithms, as well as the creators of the TwoLinkArm environment used in this study.

\begin{thebibliography}{9}

\bibitem{kalman1960contributions}
R.~E. Kalman, ``Contributions to the theory of optimal control,'' \emph{Bol. Soc. Mat. Mex.}, vol. 5, no. 2, pp. 102--119, 1960.

\bibitem{li2004iterative}
W.~Li and E.~Todorov, ``Iterative linear quadratic regulator design for nonlinear biological movement systems,'' in \emph{ICINCO}, 2004, pp. 222--229.

\bibitem{tedrake2015efficient}
R.~Tedrake, ``Efficient bipedal robots based on passive-dynamic walkers,'' \emph{Science Robotics}, vol. 350, no. 6265, pp. 269--276, 2015.

\bibitem{mason2018reusable}
S.~Mason, L.~P. Kaelbling, and T.~Lozano-P{\'e}rez, ``Reusable skills for scalable robot learning,'' \emph{arXiv preprint arXiv:1810.02863}, 2018.

\bibitem{williams2017model}
G.~Williams, P.~Drewes, B.~Goldfain, J.~M. Rehg, and E.~A. Theodorou, ``Aggressive driving with model predictive path integral control,'' in \emph{IEEE ICRA}, 2017, pp. 1433--1440.

\bibitem{howell2022do}
T.~J. Howell, C.~R. Mansley, and R.~Tedrake, ``Do differentiable simulators give better policy gradients?'' in \emph{Conference on Robot Learning}, 2022, pp. 1062--1073.

\end{thebibliography}

\appendices

\section{Mathematical Derivations}

\subsection{LQR Controller Derivation}

Starting from the infinite-horizon LQR problem, we seek to minimize the cost:
\begin{equation}
J = \int_0^\infty (\mathbf{x}^T\mathbf{Q}\mathbf{x} + \mathbf{u}^T\mathbf{R}\mathbf{u})\,dt
\end{equation}

For the linear system $\dot{\mathbf{x}} = \mathbf{A}\mathbf{x} + \mathbf{B}\mathbf{u}$, the optimal control is $\mathbf{u} = -\mathbf{K}\mathbf{x}$, where $\mathbf{K}$ satisfies the algebraic Riccati equation.

The derivation proceeds as follows:

1. Assume a quadratic value function: $V(\mathbf{x}) = \mathbf{x}^T\mathbf{P}\mathbf{x}$

2. The Hamilton-Jacobi-Bellman equation gives:
\begin{equation}
\mathbf{x}^T\mathbf{Q}\mathbf{x} + \mathbf{u}^T\mathbf{R}\mathbf{u} + \frac{\partial V}{\partial \mathbf{x}}^T(\mathbf{A}\mathbf{x} + \mathbf{B}\mathbf{u}) = 0
\end{equation}

3. Substituting the optimal control $\mathbf{u} = -\mathbf{R}^{-1}\mathbf{B}^T\frac{\partial V}{\partial \mathbf{x}}$ leads to the Riccati equation.

\subsection{iLQR Backward Pass Derivation}

The backward pass of iLQR computes the value function backwards in time. At each timestep $t$, we compute the local quadratic approximation of the value function around the current trajectory.

Given the dynamics $\mathbf{x}_{t+1} = f(\mathbf{x}_t, \mathbf{u}_t)$, the value function satisfies:
\begin{equation}
V_t(\mathbf{x}_t) = \min_{\mathbf{u}_t} \left[ l_t(\mathbf{x}_t, \mathbf{u}_t) + V_{t+1}(f(\mathbf{x}_t, \mathbf{u}_t)) \right]
\end{equation}

Linearizing around $(\mathbf{x}_t^{(k)}, \mathbf{u}_t^{(k)})$, we get:
\begin{align}
V_{t+1}(\mathbf{x}_{t+1}) &\approx V_{t+1}^{(k)} + \mathbf{v}_{\mathbf{x},t+1}^{(k)T}(\mathbf{x}_{t+1} - \mathbf{x}_{t+1}^{(k)}) \notag \\
&+ \frac{1}{2}(\mathbf{x}_{t+1} - \mathbf{x}_{t+1}^{(k)})^T\mathbf{V}_{\mathbf{xx},t+1}^{(k)}(\mathbf{x}_{t+1} - \mathbf{x}_{t+1}^{(k)})
\end{align}

Substituting the linearized dynamics and collecting terms gives the quadratic approximation of $V_t$.

\subsection{Finite-Difference Jacobian Computation}

The Jacobian matrices are computed using central differences:
\begin{equation}
\frac{\partial f_i}{\partial x_j} \approx \frac{f_i(\mathbf{x} + \epsilon\mathbf{e}_j) - f_i(\mathbf{x} - \epsilon\mathbf{e}_j)}{2\epsilon}
\end{equation}

This approximation has error $O(\epsilon^2)$ and is more accurate than forward differences. The choice of $\epsilon = 10^{-5}$ balances numerical precision with finite-precision arithmetic limitations.

\subsection{Cost Function Gradient Computation}

For the intermediate cost $l(\mathbf{x}, \mathbf{u}) = \mathbf{u}^T\mathbf{R}\mathbf{u}$, the gradients are:
\begin{align}
\frac{\partial l}{\partial \mathbf{x}} &= \mathbf{0} \\
\frac{\partial l}{\partial \mathbf{u}} &= 2\mathbf{R}\mathbf{u} \\
\frac{\partial^2 l}{\partial \mathbf{x}^2} &= \mathbf{0} \\
\frac{\partial^2 l}{\partial \mathbf{u}^2} &= 2\mathbf{R} \\
\frac{\partial^2 l}{\partial \mathbf{x}\partial \mathbf{u}} &= \mathbf{0}
\end{align}

For the terminal cost $l_f(\mathbf{x}) = (\mathbf{x} - \mathbf{x}_{goal})^T\mathbf{Q}_f(\mathbf{x} - \mathbf{x}_{goal})$, the gradients are:
\begin{align}
\frac{\partial l_f}{\partial \mathbf{x}} &= 2\mathbf{Q}_f(\mathbf{x} - \mathbf{x}_{goal}) \\
\frac{\partial^2 l_f}{\partial \mathbf{x}^2} &= 2\mathbf{Q}_f
\end{align}

\section{Implementation Details}

\subsection{Code Structure}

The implementation consists of four main modules:

1. \texttt{arm\_env.py}: Two-link arm dynamics simulation
2. \texttt{controllers.py}: LQR controller with finite-difference linearization
3. \texttt{ilqr.py}: iLQR implementation with backward and forward passes
4. \texttt{run.py}: Main execution script for experiments

\subsection{Numerical Stability Considerations}

Several measures were implemented to ensure numerical stability:

1. \textbf{Clipping}: Control inputs and states are clipped to their respective bounds
2. \textbf{Regularization}: iLQR uses Levenberg-Marquardt regularization to handle ill-conditioned Hessians
3. \textbf{Emergency Braking}: LQR includes velocity-based safety checks to prevent numerical explosions
4. \textbf{Finite-Difference Bounds}: Perturbations are kept within reasonable limits to avoid numerical issues

\subsection{Performance Optimizations}

1. \textbf{Vectorization}: All computations use NumPy vectorized operations
2. \textbf{Early Termination}: Algorithms terminate when convergence criteria are met
3. \textbf{Memory Efficiency}: Trajectories are stored efficiently without unnecessary copies
4. \textbf{Parallel Computation}: Jacobian computations could be parallelized (not implemented in this version)

\subsection{Hyperparameter Selection}

The hyperparameters were selected through a combination of theoretical analysis and empirical tuning:

\begin{itemize}
\item Finite-difference step: $\epsilon = 10^{-5}$
\item LQR state cost: $\mathbf{Q} = \text{diag}([10, 10, 1, 1])$
\item Control cost: $\mathbf{R} = \mathbf{I}_2$
\item Terminal cost: $\mathbf{Q}_f = 1000\mathbf{I}_4$
\item iLQR tolerance: $\epsilon = 10^{-3}$
\item Regularization parameters: $\lambda \in [1, 10^8]$
\end{itemize}

\section{Additional Experimental Results}

\subsection{Ablation Study}

We conducted an ablation study to understand the impact of different components:

1. \textbf{Finite-Difference Accuracy}: Increasing $\epsilon$ from $10^{-5}$ to $10^{-3}$ reduces accuracy but improves numerical stability
2. \textbf{Cost Function Design}: Higher terminal costs lead to faster convergence but more aggressive control
3. \textbf{Regularization}: Strong regularization improves stability but slows convergence
4. \textbf{Initialization}: Better initial trajectories reduce the number of iLQR iterations

\subsection{Failure Cases}

Several failure modes were observed:

1. \textbf{LQR Oscillations}: When linearization errors accumulate, LQR can produce unstable oscillatory behavior
2. \textbf{iLQR Local Optima}: Poor initialization can lead to suboptimal local minima
3. \textbf{Numerical Instability}: Large finite-difference steps can cause Jacobian computation errors
4. \textbf{Constraint Violations}: Both methods can generate infeasible trajectories if not properly constrained

\subsection{Computational Complexity Analysis}

LQR has $O(n^3)$ complexity per timestep due to Riccati equation solving, where $n$ is the state dimension. iLQR has $O(T n^3)$ complexity for the full trajectory optimization, where $T$ is the horizon length.

For real-time control, LQR is suitable for systems with $n \leq 10$, while iLQR is better suited for offline planning applications.

\end{document}

